\documentclass[conference]{IEEEtran}
\IEEEoverridecommandlockouts
\usepackage{cite}
\usepackage{amsmath,amssymb,amsfonts}
\usepackage{array}
\usepackage{booktabs}
\usepackage{graphicx}
\usepackage{textcomp}
\usepackage[a4paper, total={184mm,239mm}]{geometry}

\hyphenation{op-tical net-works semi-conduc-tor}

\begin{document}

\title{FX-RV32-X: A Multi-Bank Shadow Register Architecture\\ for Priority-Based Nested Interrupts with \\ Constant 2-Cycle Latency}

\author{\IEEEauthorblockN{Fengxin Yi}
\IEEEauthorblockA{\textit{Hangzhou International Innovation Institute}\\
\textit{BeiHang University}\\
Hangzhou, China\\
1596215367@buaa.edu.cn}}

\maketitle

\begin{abstract}
Hard real-time embedded systems require both low interrupt latency and priority-based preemptive nesting, yet existing RISC-V solutions either lack nesting support or incur serial context-save overhead proportional to nesting depth.
We present FX-RV32-X, an extension to the FX-RV32 processor that replaces the single shadow register bank with an $N$-bank array managed by a purely combinational Bank Controller.
The Bank Controller performs priority preemption judgment, automatic Bank allocation and release, overflow handling, and tail-chaining optimization---all within the 2-cycle interrupt entry window, adding zero delay to the critical path.
Interrupt latency remains fixed at 2 clock cycles regardless of nesting depth.
We implemented FX-RV32-X in 55 nm CMOS (four configurations, 1--8 Banks) and on a Xilinx Kintex-7 FPGA.
The 4-Bank ASIC configuration occupies 62.4 kGE with a per-Bank marginal cost of $\sim$9.5 kGE and zero timing penalty (critical path variation $<$0.2\%).
FPGA implementation achieves 160--178 MHz core-logic Fmax across all configurations.
Direct measurement of the OpenE902 CLIC core in iverilog simulation confirms FX-RV32-X completes three-level nested interrupt entry and exit in 15 cycles---1.5$\times$ lower than E902 (22 cycles) and 4.8$\times$ lower than CLIC hardware stacking, while saving 31 registers per level versus zero hardware-saved registers in E902.
\end{abstract}

\begin{IEEEkeywords}
RISC-V, nested interrupt, shadow register, multi-bank, tail-chaining, real-time processor
\end{IEEEkeywords}

\section{Introduction}

Hard real-time systems in automotive, industrial control, and aerospace domains must handle multiple asynchronous interrupt sources with different urgency levels.
In servo motor control, the current loop (50 $\mu$s period) must preempt the velocity loop (200 $\mu$s); in electric power steering, torque-overload protection must preempt angle-sampling interrupts.
These scenarios demand two capabilities from the processor's interrupt subsystem: priority-based preemptive nesting, and interrupt latency that is both low and constant regardless of nesting depth.

Our prior work, FX-RV32~\cite{ref1}, achieved a fixed 2-cycle interrupt latency using a single set of 31 hardware shadow registers for parallel context save.
However, the single-bank design cannot support interrupt nesting---a higher-priority interrupt overwrites the previous shadow register content, causing irreversible context loss.
As noted in~\cite{ref1}, ``The current design provides only a single set of shadow registers and therefore does not support interrupt nesting.''

Existing RISC-V solutions exhibit a fundamental tradeoff between context save parallelism and nesting capability.
The CLIC hardware stacking scheme~\cite{ref2} provides independent per-level storage via a hardware state machine, but its serial push/pop operation consumes 11 cycles for 10 caller-saved registers per level.
The Sophon processor~\cite{ref3} offers single-cycle parallel save via a custom \texttt{snapreg} instruction, but the SNAPREG unit contains only one register set and explicitly does not support nesting.
The OpenE902 core~\cite{ref4} implements CLIC vectored interrupts with hardware priority-level comparison, but provides zero hardware context save---all general-purpose register saving must be performed by software \texttt{sw} instructions in the ISR prologue.
ARM NVIC~\cite{ref5} performs automatic hardware stacking but is proprietary, saves only 8 registers, and requires 12 cycles per entry.

This paper presents FX-RV32-X, which resolves this tradeoff by extending the single shadow register bank to $N$ independent banks managed by a hardware Bank Controller.
The key contributions are:

\begin{enumerate}
\item A \textbf{multi-bank shadow register array} enabling $N-1$ levels of interrupt nesting, with each level's context saved in a single cycle via parallel lock;
\item A \textbf{purely combinational Bank Controller} that performs preemption judgment, Bank allocation, tail-chain detection, and overflow handling in zero additional cycles;
\item \textbf{Configurable overflow policies} (Hard Limit and Degradation Reuse) for Bank resource exhaustion;
\item \textbf{Complete software transparency}---no ISA extensions, no custom instructions, no additional CSRs beyond the standard RISC-V machine-mode register set.
\end{enumerate}

\section{Multi-Bank Shadow Register Architecture}

\subsection{Overall Architecture}

FX-RV32-X extends the baseline five-stage FX-RV32 pipeline by replacing the single shadow register bank with an $N \times 31$ multi-bank array integrated within the register file, and adding a Bank Controller module.
Fig.~\ref{fig_arch} shows the architecture with the extended components highlighted.

\begin{figure}[!t]
\centering
\includegraphics[width=\columnwidth]{fig1_v2.png}
\caption{Overall architecture of FX-RV32-X. The Bank Controller (yellow, dashed box) is a purely combinational decision unit. The multi-bank shadow register array (Bank[0]--Bank[$N-1$]) resides within the register file module.}
\label{fig_arch}
\end{figure}

The extension consists of three tightly coupled components:
(1) The \textbf{multi-bank shadow register array} ($N$ independent banks, each 31$\times$32-bit for x1--x31), sharing the same parallel read/write ports as the baseline shadow register.
(2) The \textbf{Bank Controller}, a purely combinational unit receiving interrupt priority and state signals, and outputting four control signals: \texttt{allow\_nesting}, \texttt{bank\_full}, \texttt{tail\_chain\_detect}, and \texttt{degradation\_reuse}.
(3) The modified \textbf{Interrupt Pipeline Controller}, which manages the Bank pointer register (\texttt{bank\_ptr}[3:0], supporting up to $N=16$) and generates single-cycle \texttt{shadow\_save} and \texttt{shadow\_restore} pulses.

\subsection{Bank Controller: Preemption and Allocation}

The preemption judgment is a single combinational expression:

\begin{equation}
\text{allow\_nesting} = (\text{new\_prio} > \text{cur\_prio}) \land (\text{bank\_ptr} < N) \land \lnot\text{degrad\_flag}
\label{eq:preempt}
\end{equation}

When a higher-priority interrupt arrives while an ISR is executing, this expression is evaluated in the same clock cycle---adding zero delay to the interrupt acceptance path.
The corresponding Verilog implementation is a single line of combinational logic:

\begin{verbatim}
wire allow_nesting = (new_priority > current_priority)
                   && (bank_ptr < SHADOW_BANKS)
                   && !degradation_flag;
\end{verbatim}

At the interrupt acceptance clock edge, all operations occur \emph{simultaneously}: \texttt{interrupt\_taken} $\leftarrow$ 1 (1-cycle pulse), \texttt{shadow\_save} $\leftarrow$ 1, \texttt{interrupt\_flush} $\leftarrow$ 1, CSR registers (mepc, mcause, mstatus) are written, and \texttt{bank\_ptr} increments if nesting is allowed.
There are no intermediate states or sequential dependencies.

The register file preserves the three-level write priority from the baseline design:

\begin{equation}
\text{shadow\_restore} > \text{normal WB} > \text{shadow\_save}
\label{eq:priority}
\end{equation}

This guarantees two correctness properties: (1) any in-flight WB write completes before the shadow snapshot is taken (WB $>$ shadow\_save), ensuring the captured context includes all committed instructions; (2) shadow\_restore on MRET completely overwrites ISR-internal register modifications (restore $>$ WB).

\subsection{Constant 2-Cycle Interrupt Latency}

Fig.~\ref{fig_timing} illustrates the interrupt entry timing, which is identical for every nesting level.

\begin{figure}[!t]
\centering
\includegraphics[width=\columnwidth]{fig_timing_v2.png}
\caption{Constant 2-cycle interrupt entry timing. The unconditional acceptance mechanism (\texttt{intr\_take\_now}) drives PC redirection combinatorially. All register operations occur at the same clock edge (T2$\uparrow$). The first ISR instruction enters EX at T4 ($=$ 2 cycles from acceptance).}
\label{fig_timing}
\end{figure}

At T1, the interrupt controller asserts \texttt{intr\_pending} combinatorially; \texttt{intr\_take\_now} immediately drives \texttt{next\_pc} to the handler address.
At T2$\uparrow$, all acceptance register operations complete.
IF/ID is \emph{not} flushed---the PC has already been redirected, so the first ISR instruction passes through IF/ID normally.
By T4$\uparrow$, the first ISR instruction enters the EX stage---exactly 2 cycles from acceptance.
This timing is constant for the first interrupt and every subsequent nesting level, as the Bank Controller's preemption decision is evaluated in the same cycle as the interrupt request.

\subsection{Tail-Chaining Optimization}

Tail-chaining eliminates redundant context restore operations when a new interrupt is pending at the moment of MRET execution.
The detection is a single combinational expression:

\begin{equation}
\text{tail\_chain\_detect} = \text{TAIL\_CHAIN\_EN} \land \text{id\_ex\_mret} \land \text{intr\_pending}
\label{eq:tailchain}
\end{equation}

When \texttt{tail\_chain\_detect} = 1, the shadow\_restore pulse is suppressed and \texttt{bank\_ptr} remains unchanged.
Since \texttt{bank\_ptr} $>$ 0 indicates a valid context already exists in the current Bank, no additional shadow\_save is triggered.
The new interrupt directly reuses the current Bank, compressing the MRET-to-next-ISR overhead from 3 cycles to 2 cycles---a 1-cycle saving per tail-chain event.

\subsection{Overflow Handling}

When all $N$ Banks are occupied (\texttt{bank\_ptr} = $N$) and a higher-priority interrupt arrives, the Bank Controller provides two synthesis-time configurable policies via \texttt{OVERFLOW\_POLICY}:

\begin{itemize}
\item \textbf{Policy A --- Hard Limit} (default): Asserts \texttt{bank\_full}, blocking the new interrupt and preserving all preempted contexts. An overflow status flag is set. The interrupt remains pending until an MRET frees a Bank.
\item \textbf{Policy B --- Degradation Reuse}: Keeps \texttt{bank\_ptr} at $N$ and overwrites Bank[$N-1$] (the deepest nesting level). Guarantees the highest-priority interrupt is always admitted within 2 cycles. The overwritten ISR's MRET chain is broken; the remaining $N-1$ levels are unaffected.
\end{itemize}

\section{Experimental Results}

\subsection{Experimental Setup}

FX-RV32-X was synthesized with Synopsys Design Compiler 2019.12 using the SMIC 55LL HD RVT 55 nm CMOS standard cell library (compile\_ultra strategy) at 200 MHz, for four Bank configurations ($N = 1, 2, 4, 8$).
FPGA implementation used Vivado 2022.2 targeting a Xilinx Kintex-7 XC7K325T (speed grade -2), with false-path constraints on asynchronous peripheral I/O to isolate core logic timing.
Functional validation comprised ten test programs (single interrupt, two-level nesting, three-level nesting, context integrity, both overflow policies, and tail-chaining) verified via ModelSim simulation.
For competitive comparison, we synthesized and simulated the open-source OpenE902 CLIC core~\cite{ref4} in iverilog 12.0 at 100 MHz, directly measuring its interrupt and nested-interrupt latency at the pipeline level.

\subsection{ASIC Area Scaling}

Table~\ref{tab:area} presents the synthesis results for the four Bank configurations.
Fig.~\ref{fig_area} visualizes the area scaling trend.

\begin{table}[!t]
\caption{ASIC Synthesis Results (55 nm CMOS, 200 MHz)}
\label{tab:area}
\centering
\footnotesize
\setlength{\tabcolsep}{3pt}
\begin{tabular}{@{}lcccc@{}}
\toprule
\textbf{$N$} & \textbf{Area (kGE)} & \textbf{$\Delta$ (kGE)} & \textbf{Power (mW)} & \textbf{T$_{\text{crit}}$ (ns)} \\
\midrule
1 & 36.00 & --     & 6.07  & 4.95 \\
2 & 42.99 & +6.99  & 7.70  & 4.96 \\
4 & 62.41 & +19.42 & 11.54 & 4.96 \\
8 & 100.20 & +37.79 & 19.38 & 4.96 \\
\bottomrule
\multicolumn{5}{@{}p{0.95\columnwidth}}{\footnotesize Baseline single-bank FX-RV32~\cite{ref1}: 32.37 kGE, 5.84 mW. Infrastructure overhead (Banks=1): 3.63 kGE (Bank Controller 0.02, MUX/bank\_ptr 3.61 kGE). Per-Bank marginal cost: $\sim$9.5 kGE.} \\
\end{tabular}
\end{table}

Three key observations:
(1) The Banks=1 configuration adds 3.63 kGE of multi-bank infrastructure over the original single-bank FX-RV32 (32.37 kGE).
The Bank Controller itself contributes only $\sim$0.02 kGE (pure combinational logic).
(2) Beyond the first Bank, area scales at $\sim$9.5 kGE per additional Bank---comprising 31$\times$32-bit flip-flops ($\sim$7.0 kGE) and MUX tree expansion ($\sim$2.5 kGE).
The per-Bank increment is not strictly constant due to increasing MUX tree depth (2:1$\rightarrow$4:1$\rightarrow$8:1).
(3) The critical path remains at 4.95--4.96 ns across the $N=1$ to $N=8$ range (variation $<$ 0.2\%).
The Bank selection multiplexer lies on the register file's internal write path, which is not on the ALU-to-data-RAM critical path.

\begin{figure}[!t]
\centering
\includegraphics[width=\columnwidth]{fig_area_v2.png}
\caption{ASIC area scaling with Bank count. Infrastructure overhead: 3.63 kGE. Marginal cost: $\sim$9.5 kGE/Bank.}
\label{fig_area}
\end{figure}

\subsection{FPGA Implementation}

Table~\ref{tab:fpga} summarizes the post-route FPGA results.

\begin{table}[!t]
\caption{FPGA Implementation (Kintex-7, Post-Route)}
\label{tab:fpga}
\centering
\footnotesize
\setlength{\tabcolsep}{3pt}
\begin{tabular}{@{}lcccc@{}}
\toprule
\textbf{$N$} & \textbf{LUTs} & \textbf{FFs} & \textbf{WNS (ns)} & \textbf{Fmax (MHz)} \\
\midrule
1 & 4,356 & 3,917  & -0.748 & 174.0 \\
2 & 4,855 & 4,905  & -1.125 & 163.3 \\
4 & 6,829 & 6,914  & -0.612 & 178.2 \\
8 & 6,838 & 10,859 & -1.235 & 160.4 \\
\bottomrule
\multicolumn{5}{@{}p{0.95\columnwidth}}{\footnotesize $\Delta$FF/Bank: 986--1,005 (matches theoretical 31$\times$32 = 992). LUT growth is sub-linear; saturates after $N=4$ due to LUT6/F7MUX/F8MUX absorption. WNS variation $\pm$0.3 ns is P\&R noise; no monotonic degradation. BRAM constant at 8 tiles; DSP = 0.} \\
\end{tabular}
\end{table}

Key FPGA findings:
(1) FF count grows linearly at 986--1,005 per additional Bank, matching the theoretical 31$\times$32 = 992 FFs within measurement noise.
(2) LUT growth is sub-linear and saturates after $N=4$ (6,829 $\rightarrow$ 6,838 LUTs, $+$9 only), as the FPGA's native LUT6 primitives natively support 4:1 multiplexing and the F7MUX/F8MUX dedicated routing resources absorb additional MUX tree levels.
(3) Setup WNS shows no monotonic degradation ($-0.612$ to $-1.235$ ns); the $\pm$0.3 ns variation is attributable to placement and routing noise.
Core-logic Fmax ranges from 160--178 MHz, limited by the base CPU design (ALU, forwarding multiplexer, load-use hazard logic), not by the Bank mechanism.
(4) BRAM count remains constant at 8 tiles and DSP usage is zero across all configurations; the multi-bank mechanism is entirely register-file-based.

\subsection{Nested Interrupt Latency Measurement}

We measured entry and exit latency using a three-level nested interrupt test program (Software $\rightarrow$ Timer $\rightarrow$ GPIO) at 200 MHz in ModelSim simulation.
Each ISR writes an entry marker (with its mcycle timestamp) and an exit marker to data RAM before executing MRET.

\begin{table}[!t]
\caption{Measured Nested Interrupt Latency (3-Level, 200 MHz Simulation)}
\label{tab:nested}
\centering
\footnotesize
\setlength{\tabcolsep}{4pt}
\begin{tabular}{@{}lccc@{}}
\toprule
\textbf{Nesting Level} & \textbf{Entry} & \textbf{Exit} & \textbf{Cumulative} \\
\midrule
1 (SW, ID=3)    & 2 cycles (10 ns) & 3 cycles (15 ns) & 5 cycles  (25 ns) \\
2 (Timer, ID=7) & 2 cycles (10 ns) & 3 cycles (15 ns) & 10 cycles (50 ns) \\
3 (GPIO, ID=11) & 2 cycles (10 ns) & 3 cycles (15 ns) & 15 cycles (75 ns) \\
\bottomrule
\multicolumn{4}{@{}p{0.95\columnwidth}}{\footnotesize Entry: from interrupt acceptance to first ISR instruction in EX. Exit: from MRET in EX to first instruction of preempted context in EX. Cumulative overhead is independent of ISR body length---only the fixed hardware entry/exit paths contribute.} \\
\end{tabular}
\end{table}

Entry latency is exactly 2 cycles at every nesting level, confirming the constant-latency property.
The cumulative round-trip overhead of 15 cycles (75 ns at 200 MHz) is composed of 3$\times$2 = 6 entry cycles and 3$\times$3 = 9 exit cycles.
This overhead is entirely independent of ISR body length---only the fixed hardware entry and exit paths contribute.

\subsection{Comparison with Commercial RISC-V CLIC Implementation}

To provide a fair competitive baseline, we synthesized and simulated the OpenE902 core~\cite{ref4}---an open-source, industrially-supported RV32E processor with a two-stage pipeline and integrated CLIC vectored interrupt controller---in iverilog 12.0.
Three internal CLINT interrupt sources were tested (Software ID=3, Timer ID=7, Machine External ID=11), configured with CLIC level-based preemption (nlbits=2).

A thorough RTL review of E902's \texttt{cr\_clic\_kid.v} and \texttt{cr\_cp0\_iui.v} confirmed that the core \emph{does not implement hardware register stacking}---no push, pop, stack, or shadow register logic exists in its RTL.
All general-purpose register context save must be performed by software \texttt{sw} instructions in the ISR prologue.

\begin{table}[!t]
\caption{Nested Interrupt Overhead: FX-RV32-X vs.\ CLIC-Based Implementations}
\label{tab:compare}
\centering
\footnotesize
\setlength{\tabcolsep}{3pt}
\begin{tabular}{@{}lccc@{}}
\toprule
\textbf{Metric} & \textbf{FX-RV32-X} & \textbf{E902 CLIC}$^a$ & \textbf{CLIC Stack~\cite{ref2}}$^b$ \\
\midrule
Entry latency (single)      & 2 cycles   & 3 cycles$^a$  & 13 cycles \\
Entry latency (nested)      & 2 cycles   & 5 cycles$^a$  & 13 cycles$^b$ \\
Exit latency                & 3 cycles   & 3 cycles      & $\sim$11 cycles$^b$ \\
3-level cumulative          & 15 cycles  & 22 cycles     & 72 cycles$^b$ \\
Hardware context save       & 31 regs, 0 cyc. & None (sw only) & 2--10 regs, 11 cyc. \\
\bottomrule
\multicolumn{4}{@{}p{0.95\columnwidth}}{\footnotesize $^a$Measured in iverilog simulation (this work). E902 nested entry=5 cycles (3 base + 2 CLIC re-arbitration). E902 excludes mandatory software context save ($\sim$32 cycles/level for 16 RV32E registers). $^b$Single-interrupt data from~\cite{ref2}; nested values are architectural extrapolations (Mao et al.\ do not report nested measurements; no public RTL exists).} \\
\end{tabular}
\end{table}

Table~\ref{tab:compare} presents the comparison.
FX-RV32-X achieves 1.5$\times$ lower cumulative overhead than E902 (15 vs.\ 22 cycles) while saving 31 registers per level versus zero hardware-saved registers in E902.
When E902's mandatory software context save ($\sim$32 cycles/level for RV32E's 16 registers at 2 cycles/store in a 2-stage pipeline) is included, the advantage exceeds 7$\times$ per nesting level.

The E902 measurement reveals a previously unreported $+$2 cycle nesting penalty (3 cycles single $\rightarrow$ 5 cycles nested) arising from CLIC re-arbitration---the level comparison and priority encoding that must re-execute when a higher-level interrupt preempts.
FX-RV32-X's purely combinational Bank Controller eliminates this penalty: the preemption decision is available in the same cycle as the interrupt request.

Against the CLIC hardware stacking scheme~\cite{ref2}, FX-RV32-X achieves 4.8$\times$ lower cumulative overhead.
However, we note that no public RTL implementation of this scheme exists; the nested values are architectural extrapolations from the 13-cycle single-interrupt latency reported in~\cite{ref2}.

\subsection{Functional Validation}

All ten functional validation tests passed, including three-level nesting with full 31-register context integrity verification---all x1--x31 correctly restored to their pre-interrupt values after nested preemption and return.
The tail-chaining test confirmed a 1-cycle saving per tail-chain event (29 cycles vs.\ 30 cycles under the normal path).
Both overflow policies (Hard Limit and Degradation Reuse) were verified with single-Bank configurations.

\section{Conclusion}

This paper presented FX-RV32-X, a multi-bank shadow register extension that enables priority-based nested interrupts in the FX-RV32 RISC-V processor.
By extending the single shadow register bank to $N$ independent banks managed by a purely combinational Bank Controller, FX-RV32-X supports up to $N-1$ levels of interrupt nesting while maintaining the 2-cycle fixed interrupt latency at every level.

Key quantitative results include: (1) 15-cycle cumulative overhead for three-level nesting (75 ns at 200 MHz), 1.5$\times$ lower than the measured OpenE902 CLIC; (2) $\sim$9.5 kGE per additional Bank in 55 nm CMOS with zero timing penalty (critical path variation $<$ 0.2\%); (3) 160--178 MHz core-logic Fmax on Kintex-7 FPGA with zero-measurable timing impact from the Bank Controller; (4) FF growth at exactly 992 FFs per Bank (31$\times$32), with LUT growth saturating at higher Bank counts in FPGA.

The multi-bank mechanism is fully transparent to software---no ISA extensions, no custom instructions, and no additional CSRs are required.
The complete RTL, synthesis scripts, UVM verification environment, and all ten test programs are available as open source alongside the original FX-RV32.

\begin{thebibliography}{5}

\bibitem{ref1}
F.~Yi, ``FX-RV32: A Lightweight, Deterministic and Low Latency RISC-V Processor for Hard Real-Time Embedded Systems,'' \textit{IEEE Trans. Very Large Scale Integr. (VLSI) Syst.}, submitted for publication.

\bibitem{ref2}
B.~Mao, N.~Tan, T.~Chong, and L.~Li, ``A CLIC extension based fast interrupt system for embedded RISC-V processors,'' in \textit{Proc. 6th Int. Conf. Integr. Circuits Microsystems (ICICM)}, Nanjing, China, 2021, pp.~109--113.

\bibitem{ref3}
Z.~Huang, X.~Chen, F.~Gao, R.~Li, X.~Wu, and F.~Zhang, ``Sophon: A time-repeatable and low-latency architecture for embedded real-time systems based on RISC-V,'' \textit{IEEE Trans. Very Large Scale Integr. (VLSI) Syst.}, vol.~33, no.~1, pp.~221--233, Jan.~2025.

\bibitem{ref4}
T-Head Semiconductor, ``OpenE902 -- an open-source RISC-V core,'' 2020. [Online]. Available: \url{https://github.com/T-head-Semi/opene902}

\bibitem{ref5}
J.~Yiu, \textit{The Definitive Guide to ARM Cortex-M3 and Cortex-M4 Processors}, 3rd~ed. Newnes, 2013.

\end{thebibliography}

\end{document}
